% ==================================================================
%  WEEK 9
% ==================================================================
\chapter{Hierarchical Edge-Fog-Cloud Computing and Cloud Federation}

\section{Introduction and Week Overview}
Week 9 explores the continuum between centralized hyperscale data centers and decentralized edge computing. Centralized clouds struggle to process the massive volume, velocity, and variety of data generated by billions of Internet of Things (IoT) devices with low latency. This module covers:
\begin{enumerate}[leftmargin=*]
  \item \textbf{The Edge-Fog-Cloud Continuum:} Cisco's Fog Computing architectural paradigm, multi-tier hierarchical delegation, and latency-sensitive control loops.
  \item \textbf{Empirical Performance Modeling with iFogSim:} Network bandwidth conservation, execution cost reduction, latency bounds, and healthcare case studies.
  \item \textbf{Resource Management and Service Placement:} Formal taxonomies of resource management (Data Flow, Control, Tenancy), vertical vs horizontal scalability, and multi-objective task placement.
  \item \textbf{Cloud Federation and Inter-Cloud Governance:} Overcoming vendor lock-in through Broker Architectures, Peer-to-Peer (P2P) peering, and cross-provider SLA enforcement.
\end{enumerate}

\section{The Edge-Fog-Cloud Computing Hierarchy}
\subsection{The Need for Decentralization}
While hyperscale cloud computing provides elastic computing capacity and virtually infinite storage, routing all raw IoT telemetry to centralized data centers introduces critical bottlenecks:
\begin{itemize}[leftmargin=*]
  \item \textbf{Excessive WAN Latency:} Multi-hop packet routing across public Internet backbones introduces unpredictable delays ($100\text{ ms}$ to $300\text{ ms}$), violating real-time cyber-physical feedback loops.
  \item \textbf{Core Network Congestion:} Streaming millions of high-definition video feeds or high-frequency vibrational sensor telemetry overloads Internet Service Provider (ISP) backbones.
  \item \textbf{Data Privacy and Sovereignty:} Raw audio, video, and biometric streams transmitted across national borders pose regulatory and privacy risks.
\end{itemize}

\subsection{Architectural Division of Responsibilities}
Introduced by Cisco Systems and formalized in IEEE/ACM distributed systems research, the \textbf{Edge--Fog--Cloud} hierarchy divides responsibilities across three distinct tiers:

\begin{figure}[H]
\centering
\begin{tikzpicture}[font=\small]
  % Cloud Tier
  \node[dark, minimum width=12cm, minimum height=1.1cm] (cloud) {\textbf{Tier 3: Centralized Hyperscale Cloud}\\Global orchestration, deep machine learning training, petabyte warehousing ($> 100\text{ ms}$ latency)};
  
  % Fog Tier
  \node[mid, below=1cm of cloud, minimum width=12cm, minimum height=1.1cm] (fog) {\textbf{Tier 2: Fog Computing Layer (Area Gateways, Roadside Micro Data Centers)}\\Regional data aggregation, anomaly detection, localized cross-device coordination ($10\text{ ms} - 50\text{ ms}$ latency)};
  
  % Edge Tier
  \node[gold, below=1cm of fog, minimum width=12cm, minimum height=1.1cm] (edge) {\textbf{Tier 1: Edge / IoT Device Layer (Sensors, Actuators, Smart Cameras, PLCs)}\\Immediate local filtering, real-time closed-loop control, raw telemetry pruning ($< 10\text{ ms}$ latency)};

  \draw[arr, <->] (cloud) -- node[right, font=\scriptsize] {Summarized Insights / Batch Analytics} (fog);
  \draw[arr, <->] (fog) -- node[right, font=\scriptsize] {Filtered Alerts / Local Telemetry} (edge);
\end{tikzpicture}
\caption{Three-Tier Edge-Fog-Cloud Hierarchical Architecture.}
\end{figure}

\begin{center}
\begin{tabularx}{\textwidth}{>{\bfseries}l X X X}
\toprule
System Dimension & Edge Layer & Fog Layer & Cloud Layer\\
\midrule
Response Latency & Ultra-low ($1\text{ ms} - 10\text{ ms}$) & Low ($10\text{ ms} - 50\text{ ms}$) & Moderate to High ($100\text{ ms} - 500\text{ ms}$)\\
Data Scope & Transient, sensor-local & Regional, facility-wide & Global, enterprise-wide\\
Storage Duration & Seconds (transient buffer) & Hours to days & Years (permanent historical archival)\\
Compute Hardware & Microcontrollers, ASICs, Raspberry Pi & Micro-servers, edge gateways, high-end routers & Multi-socket server blades, GPU clusters\\
Data Volume Flow & Ingests $100\%$ of raw telemetry & Filters/aggregates to $\approx 10\%$ of raw data & Receives $< 1\%$ (compact statistical summaries)\\
Network Uplink & Local Wi-Fi, BLE, ZigBee & Regional Ethernet, 5G NR, Fiber & Transoceanic Internet WAN fiber\\
\bottomrule
\end{tabularx}
\end{center}

\subsection{Empirical Simulation Analysis: The iFogSim Framework}
The lecture features research co-authored by Prof. Soumya K. Ghosh and Prof. Rajkumar Buyya utilizing \textbf{iFogSim}, an open-source Java simulation toolkit for modeling resource management in IoT, Edge, and Fog environments:

\subsubsection{Healthcare Case Study: Smart Ambient Assisted Living}
A patient monitoring application measures ECG, blood pressure, and pulse telemetry:
\begin{itemize}[leftmargin=*]
  \item \textbf{Configuration A (Cloud-Only):} Every raw biometric pulse packet is forwarded directly to AWS cloud servers.
  \item \textbf{Configuration B (Cloud-Fog-Edge):} Wearable edge sensors execute lightweight feature extraction. An Area Gateway (Fog Node) runs automated arrhythmia detection algorithms locally. Only abnormal readings trigger an alert that contacts the cloud confirmatory module.
\end{itemize}

\subsubsection{Performance Findings}
\begin{enumerate}[leftmargin=*]
  \item \textbf{Network Bandwidth Consumption:} In the Fog configuration, network usage decreases by over $90\%$ because raw sinusoidal signals are evaluated locally; only verified anomaly events traverse the wide-area network.
  \item \textbf{Operational Execution Cost:} Cloud compute hours are strictly pay-per-use. Processing raw data on customer-owned edge gateways slashes recurring monthly cloud billing fees.
  \item \textbf{Control Loop Latency:} In the Cloud-only model, closed-loop actuation delays fluctuated between $1500\text{ ms}$ and $4000\text{ ms}$ due to network congestion. In the Fog architecture, control loop latency remained deterministic and fixed below $50\text{ ms}$.
\end{enumerate}

\section{Resource Management and Scalability Models}
\subsection{Resource Management Taxonomies}
Following the ACM survey by Hong and Varghese highlighted in the lecture, resource management in Fog/Edge computing is partitioned along three axes:
\begin{itemize}[leftmargin=*]
  \item \textbf{Data Flow:} Movement of workloads and telemetry. Workloads may flow upstream (from sensor devices to edge/cloud) or downstream (offloading pre-trained AI weights from cloud to edge inferencing engines).
  \item \textbf{Control Hierarchy:}
    \begin{itemize}
      \item \emph{Centralized Control:} A master coordinator governs scheduling across all edge devices.
      \item \emph{Distributed / Autonomous Control:} Edge and fog nodes negotiate tasks peer-to-peer using localized game theory or consensus algorithms.
    \end{itemize}
  \item \textbf{Tenancy:} Single-tenant dedicated hardware vs multi-tenant virtualized containers running concurrent multi-user tasks.
\end{itemize}

\subsection{Vertical vs Horizontal Scalability}
When a computing node experiences resource exhaustion:
\begin{defbox}{Vertical Scalability (Scale-Up) vs Horizontal Scalability (Scale-Out)}
\textbf{Vertical Scalability (Scale-Up):} Increasing the physical or virtual capacity of an existing node by assigning additional CPU cores, expanding system RAM, or upgrading disk I/O throughput within the same machine.\\
\textbf{Horizontal Scalability (Scale-Out):} Adding new independent nodes or servers into the computing pool and distributing incoming workload requests across the expanded cluster via a load balancer.
\end{defbox}

\section{Cloud Federation and Inter-Cloud Governance}
\subsection{Motivation for Cloud Federation}
Relying on a single Cloud Service Provider (CSP) exposes enterprises to catastrophic risks:
\begin{itemize}[leftmargin=*]
  \item \textbf{Vendor Lock-In:} Proprietary APIs, incompatible storage schemas, and exorbitant egress bandwidth fees trap enterprise applications.
  \item \textbf{Geographic Blackouts:} Outages in a single provider's regional data center can incapacitate global business operations.
  \item \textbf{Capacity Exhaustion:} Sudden global demand spikes can surpass a single provider's available unallocated capacity.
\end{itemize}

\begin{defbox}{Cloud Federation}
\textbf{Cloud Federation} is the seamless aggregation and dynamic coordination of computing, storage, and networking resources from multiple autonomous, heterogeneous Cloud Service Providers to satisfy consumer requirements under a unified SLA.
\end{defbox}

\subsection{Federation Architectural Topologies}
\begin{figure}[H]
\centering
\begin{tikzpicture}[font=\small]
  % Broker Model
  \node[dark, minimum width=2.4cm] (c1) {Consumer};
  \node[gold, right=1.5cm of c1, minimum width=2.8cm, minimum height=1cm] (broker) {\textbf{Federation Broker}\\(SLA Matchmaker)};
  \node[mid, right=1.5cm of broker, yshift=1cm, minimum width=2.2cm] (csp1) {CSP A (AWS)};
  \node[mid, right=1.5cm of broker, minimum width=2.2cm] (csp2) {CSP B (Azure)};
  \node[mid, right=1.5cm of broker, yshift=-1cm, minimum width=2.2cm] (csp3) {CSP C (GCP)};

  \draw[arr, <->] (c1) -- (broker);
  \draw[arr, <->] (broker) -- (csp1);
  \draw[arr, <->] (broker) -- (csp2);
  \draw[arr, <->] (broker) -- (csp3);

  \node[below=1.5cm of c1, font=\bfseries] {(a) Broker-Based Federation};

  % P2P Model
  \node[mid, right=6.5cm of c1, yshift=0.8cm, minimum width=2.2cm] (p1) {CSP 1};
  \node[mid, right=2cm of p1, minimum width=2.2cm] (p2) {CSP 2};
  \node[mid, below=1.2cm of $(p1)!0.5!(p2)$, minimum width=2.2cm] (p3) {CSP 3};

  \draw[arr, <->] (p1) -- node[above, font=\scriptsize] {Peering SLA} (p2);
  \draw[arr, <->] (p2) -- (p3);
  \draw[arr, <->] (p3) -- (p1);

  \node[below=0.6cm of p3, font=\bfseries] {(b) Peer-to-Peer (P2P) Federation};
\end{tikzpicture}
\caption{Architectural Models of Cloud Federation.}
\end{figure}

\begin{enumerate}[leftmargin=*]
  \item \textbf{Broker Architecture (Centralized Federation):} An independent entity (the Cloud Federation Broker) acts as a neutral intermediary between cloud consumers and multiple CSPs:
    \begin{itemize}
      \item Discovers available capacity and catalogs service pricing across diverse providers.
      \item Evaluates consumer QoS and latency constraints, selecting the optimal provider.
      \item Negotiates multi-provider contractual SLAs and orchestrates cross-cloud billing.
    \end{itemize}
  \item \textbf{Peer-to-Peer (P2P) Federation:} Cloud service providers negotiate direct peering agreements without a central broker. When Provider A experiences localized capacity exhaustion, it transparently subcontracts surplus workload to Provider B under pre-agreed bilateral cross-charging tariffs.
\end{enumerate}

\section{Comprehensive Solved Assignment Problems (Week 9)}

\begin{solved}{Multi-Tier Placement for Industrial Anomaly Detection (Assignment 9, Q1)}
\Given A factory requires machine-level anomaly detection within milliseconds and plant-wide data aggregation before sending summaries to the cloud. What is the optimal placement?
\Steps
\begin{enumerate}[leftmargin=*]
  \item Millisecond anomaly detection on individual machines requires processing at the \textbf{Edge}.
  \item Plant-wide aggregation across multiple machines requires intermediate coordination at \textbf{Fog} nodes.
  \item Final long-term historical reporting is handled in the \textbf{Cloud}.
\end{enumerate}
\Answer \textbf{Option B: Machine-level processing at the edge and plant-level processing at fog nodes}.
\end{solved}

\begin{solved}{Smart City Multi-Tier Video Analytics Architecture (Assignment 9, Q2)}
\Given Object detection occurs near cameras, regional congestion patterns are analyzed at roadside micro-datacenters, and long-term city trends are processed in the cloud. Which architecture represents this system?
\Steps
\begin{enumerate}[leftmargin=*]
  \item Processing near cameras $\to$ Edge.
  \item Roadside micro-datacenters $\to$ Fog.
  \item Long-term trends in the central facility $\to$ Cloud.
  \item The complete multi-tiered paradigm is \textbf{Edge--Fog--Cloud}.
\end{enumerate}
\Answer \textbf{Option C: Edge--Fog--Cloud architecture}.
\end{solved}

\begin{solved}{Raw Data Volume Transmitted to Central Infrastructure (Assignment 9, Q3)}
\Given Rank the architectural strategies by raw data volume transmitted toward the central cloud, from lowest to highest.
\Steps
\begin{enumerate}[leftmargin=*]
  \item \textbf{Edge processing} filters and prunes telemetry immediately at the sensor, transmitting the least volume.
  \item \textbf{Fog processing} aggregates data across a local area, transmitting moderate volume.
  \item \textbf{Cloud-only processing} streams all raw uncompressed data over the WAN, transmitting the highest volume.
\end{enumerate}
\Answer \textbf{Option A: Edge processing $\to$ Fog processing $\to$ Cloud-only processing}.
\end{solved}

\begin{solved}{Scaling Mechanism Classification (Assignment 9, Q4)}
\Given A fog node running low on memory is provisioned additional RAM and CPU cores on that same physical machine. Which scalability approach is this?
\Steps
\begin{enumerate}[leftmargin=*]
  \item Upgrading the CPU and RAM capacity of an existing machine without adding new nodes is \textbf{Vertical Scalability} (Scale-Up).
\end{enumerate}
\Answer \textbf{Option D: Vertical scalability}.
\end{solved}

\begin{solved}{Division of Responsibilities in Edge-Fog-Cloud (Assignment 9, Q5)}
\Given Which statement best describes the division of responsibilities in a hierarchical Edge-Fog-Cloud system?
\Steps
\begin{enumerate}[leftmargin=*]
  \item Edge executes immediate local actions.
  \item Fog coordinates regional resources and aggregates telemetry.
  \item Cloud performs global optimization and long-term strategic analytics.
\end{enumerate}
\Answer \textbf{Option B: Edge handles immediate local actions, Fog coordinates regional resources, and Cloud performs global orchestration}.
\end{solved}

\begin{solved}{Split Architecture for Video Analytics (Assignment 9, Q6)}
\Given A video analytics application requires first-stage inference within 10 ms of the camera, while heavy model retraining requires massive computing power.
\Steps
\begin{enumerate}[leftmargin=*]
  \item Sub-10 ms inference must run at the Edge/Fog close to the camera.
  \item Heavy model training runs in the centralized Cloud.
\end{enumerate}
\Answer \textbf{Option C: First-stage inference on an edge/fog node and model training in the cloud}.
\end{solved}

\begin{solved}{Federation Broker Responsibilities (Assignment 9, Q7)}
\Given What is the primary role of a Cloud Federation Broker?
\Steps
\begin{enumerate}[leftmargin=*]
  \item The Federation Broker mediates between consumers and multiple CSPs, discovering available resources, matching SLA terms, and routing workloads across cloud providers.
\end{enumerate}
\Answer \textbf{Option A: Discovers, negotiates, and coordinates services across multiple independent cloud providers}.
\end{solved}

\begin{solved}{Limitation of Edge Devices in Heavy Analytics (Assignment 9, Q8)}
\Given Why can't edge devices replace centralized cloud data centers for large-scale historical model training?
\Steps
\begin{enumerate}[leftmargin=*]
  \item Edge devices have constrained CPU, GPU, memory, and battery capacity, making them incapable of processing multi-terabyte datasets required for model training.
\end{enumerate}
\Answer \textbf{Option B: Edge devices possess limited computational power, memory, and storage capacity}.
\end{solved}

\begin{solved}{Peer-to-Peer vs Broker Federation (Assignment 9, Q9)}
\Given In a Peer-to-Peer (P2P) cloud federation model, how are resources negotiated?
\Steps
\begin{enumerate}[leftmargin=*]
  \item Participating cloud providers establish direct, bilateral peering agreements to share surplus capacity without relying on a centralized intermediary broker.
\end{enumerate}
\Answer \textbf{Option C: Directly between participating cloud providers through bilateral peering agreements}.
\end{solved}

\begin{solved}{Bandwidth Conservation via Fog Preprocessing (Assignment 9, Q10)}
\Given How does fog preprocessing conserve WAN network bandwidth?
\Steps
\begin{enumerate}[leftmargin=*]
  \item By filtering, aggregating, and compressing raw sensor streams locally, preventing redundant telemetry from being transmitted across the core Internet backbone.
\end{enumerate}
\Answer \textbf{Option A: By filtering and aggregating telemetry locally before wide-area transmission}.
\end{solved}
